% udt — example of the report mode
\documentclass[a4paper, polish, 11pt]{scrartcl}
\usepackage{udt-report}

\course{Laboratorium drgań i zjawisk falowych}
\exercise{F-00}
\title{Wahadło matematyczne}
\author{Jan Kowalski}
\studentid{000000}
\supervisor{dr Jan Nowak}
\experimentdate{01.10.2026}
\date{08.10.2026}

\begin{document}
\maketitle

\begin{preface}
Celem ćwiczenia jest wyznaczenie przyspieszenia ziemskiego z okresu
drgań wahadła matematycznego.
\end{preface}

\begincolumns

\section{Wstęp}\label{sec:wstep}
Dla małych wychyleń okres drgań wahadła o długości $l$ wynosi
\begin{equation}
	T = 2\pi\sqrt{\frac{l}{g}}
	\label{eq:okres}
\end{equation}
skąd po przekształceniu \cref{eq:okres} otrzymujemy $g = 4\pi^2 l / T^2$.

\section{Pomiary}\label{sec:pomiary}
Wyniki zebrano w \cref{tab:pomiary}, a zależność $T^2(l)$ pokazano
na \cref{fig:wykres}.

\begin{table}
	\centering
	\caption{Okres drgań dla różnych długości wahadła}
	\label{tab:pomiary}
	\begin{tabular}{S[table-format=1.2] S[table-format=1.3]}
		\toprule
		{$l$ [\si{\metre}]} & {$T$ [\si{\second}]} \\
		\midrule
		0.50 & 1.419 \\
		0.75 & 1.738 \\
		1.00 & 2.006 \\
		\bottomrule
	\end{tabular}
\end{table}

\begin{figure}
	\centering
	\begin{tikzpicture}
		\begin{axis}[width=\columnwidth, xlabel={$l$ [\si{\metre}]}, ylabel={$T^2$ [\si{\second\squared}]}]
			\addplot[only marks] coordinates {(0.50,2.014) (0.75,3.021) (1.00,4.024)};
			\addplot[domain=0.45:1.05] {4.025*x};
		\end{axis}
	\end{tikzpicture}
	\caption{Zależność kwadratu okresu od długości wahadła}
	\label{fig:wykres}
\end{figure}

\section{Wnioski}
Z dopasowania prostej (\nameref{sec:pomiary}) otrzymano
$g = \qty{9.81(4)}{\metre\per\second\squared}$.

\end{document}
